\section{Hash Security of the Root Tag}
\label{sec:cmtds3}

The tag-level results of this section share one spine. The Root tag
Candidate Sequence lemma (\Cref{lem:root-candidate-seq}) collects the
root tag inputs reached by a keyless game into a chronological candidate
sequence, and the Root Tag Hash Bound (\Cref{lem:root-tag-hash}) turns a
collision, preimage, or second-preimage among those candidates into the
matching adaptive bound of \Cref{app:lemmas}, controlling the chance that
distinct inputs land on the same $\tauroot$-bit root tag projection. The
four random oracle theorems---CMTDs-4 (\Cref{thm:cmtds4-ro}), collision
(\Cref{thm:coll-ro}), preimage (\Cref{thm:epre-ro}), and
always second-preimage (\Cref{thm:asec-ro})---are then short
instantiations, differing only in which case of the bound they invoke and in
whether a leaf tag collision term is needed when colliding inputs may use
different ciphertext bodies. The context-discoverability proof of
\Cref{thm:cdy-full} uses the same fixed-target root-preimage case,
together with hidden-component key-test accounting.
\Cref{app:flat-sponge-lift} transfers each instance to the public
permutation setting, \Cref{sec:headline} records the concrete bounds, and
\Cref{cor:committing-combined} forwards the remaining \cite{BH22}
committing notions to the same root-collision bound.

The instantiations rest on two structural separation facts, recorded
first: distinct contexts $(K, N, A)$ reach distinct final root
transcripts regardless of the leaf tags they aggregate (Opening Triple
Separation, \Cref{cor:triple-sep}), and---absent a collision among those
leaf tags---so do distinct full inputs (Final-Root Input Separation,
\Cref{prop:final-root-sep}).

\begin{proposition}[Opening Triple Separation]
\label{cor:triple-sep}
If two \TW{} root computations have distinct $(K, N, A)$ triples, then
their bridged final root transcript inputs under $\flatmap(\cdot)$ are
distinct, regardless of whether each computation is an encryption or
a decryption check on a common ciphertext.
\end{proposition}

\begin{proof}
Suppose for contradiction that $\flatmap(R_i) = \flatmap(R_j)$. Since the
bridge $\flatmap(\cdot)$ is injective on well-formed transcript prefixes
(\Cref{lem:bridge-decode}), the native flattened inputs of the two
final-root prefixes are equal. Then \Cref{lem:transcript-inj}
recovers the full duplex-call sequence and reads off the pair
$(\sigma,\sigph)$ for each call. The seven $4$-bit phase suffixes are
pairwise distinct, so the \siginit, $\sigma_{\mathsf{ad}}^\pm$,
$\sigma_{\mathsf{msg}}^\pm$, and (when present) $\sigma_{\mathsf{agg}}^\pm$
calls in the recovered sequence are unambiguously identified. The
aggregation phase is present exactly when $n \geq 1$, and
$n = \leaves(C)$ is determined by the common ciphertext, so the two
recovered sequences either both carry an aggregation phase or both omit
it; either way the recovery of $(K, N, A)$ below uses only the
\siginit and AD calls, which precede any message or aggregation block.

Equality of the recovered sequences forces equality of the \siginit
call, whose $\sigma$ is $\prfxroot \concat K \concat N$ with
fixed-width $K$ and $N$, hence $K_i = K_j$ and $N_i = N_j$. It also
forces equality of the recovered AD subsequences. By the AD bullet of
\Cref{lem:transcript-inj} the associated data phase contributes a
$\sigadlast$-tagged block exactly when the associated data is nonempty,
so the recovered sequences agree on the presence of the AD phase. If
exactly one of $A_i, A_j$ were empty the recovered sequences would
differ in the presence of a $\sigadlast$ call, contradicting their
equality; hence either both are empty, giving $A_i = A_j = \varepsilon$,
or both are nonempty, in which case the phase-recovery sub-claim of
\Cref{lem:transcript-inj} inverts the $\rho$-byte-block partition map on
the AD byte string and equal AD $\sigma$-block sequences imply
$A_i = A_j$. In all cases $(K_i, N_i, A_i) = (K_j, N_j, A_j)$,
contradicting the hypothesis of distinct triples.
\end{proof}

Triple Separation handles distinct contexts; when inputs share a context
but differ in their messages, distinctness of the final root transcripts
needs the absence of a leaf tag collision, recorded next.

\begin{proposition}[Final-Root Input Separation]
\label{prop:final-root-sep}
Let $\badleaf^\star$ be the event that two distinct construction-generated
final leaf transcripts receive the same $\tauleaf$-bit $\RF$ projection.
On $\neg\badleaf^\star$, if two \TW{} construction computations reach
final root transcripts with $\flatmap(R) = \flatmap(R')$, then they share
the same context $(K, N, A)$ and the same full ciphertext $C$.
Equivalently, on $\neg\badleaf^\star$ computations with distinct
$(K, N, A, C)$ reach distinct bridged final root transcripts.
\end{proposition}

\begin{proof}
By bridge injectivity (\Cref{lem:bridge-decode}),
$\flatmap(R) = \flatmap(R')$ gives equal native flattened final-root
inputs. \Cref{lem:transcript-inj} then recovers the same root
initialization $\prfxroot \concat K \concat N$---hence the same $K$ and
$N$---the same associated data blocks---hence the same $A$---the same root
ciphertext chunk, and the same leaf count and ordered leaf tag sequence.
Under $\neg\badleaf^\star$, Leaf-Transcript Recovery
(\Cref{lem:leaf-recovery}) identifies the equal leaf tag sequences with
the same final leaf transcripts, hence the same leaf ciphertext bodies;
with the common root ciphertext chunk this gives the same full $C$.
\end{proof}

\begin{lemma}[Root tag Candidate Sequence]
\label{lem:root-candidate-seq}
Consider a random oracle game of \Cref{sec:ro-model} in which the
adversary $\mathcal{B}$ makes direct $\ROflat$ queries and the challenger
runs $m \geq 1$ deterministic
encryption or decryption checks, the $i$-th reaching a final root transcript $R_i$
whose root tag is read by a single $\ROflat$ lookup at $\flatmap(R_i)$.
Every execution then admits a chronological sequence of root tag inputs
such that
\begin{enumerate}
\item each distinct direct query whose input $Y$ decodes under
  $\KeyedTWOut$ to a root tag output-point class---$\Rmsg^+$ (the final
  root message block, which carries the root tag only on the no-leaf path)
  or $\Rtag$ (the aggregated root tag)---is appended as a candidate before
  its $\ROflat$ lookup, and contributes at most one candidate; multi-chunk
  computations also reach $\Rmsg^+$ with zero requested output before
  aggregation, so this predicate over-approximates the root tag candidate
  set, which only loosens the count;
\item if the challenger reaches its checks, each $\flatmap(R_i)$ is in
  the sequence, appended---if not already present---immediately before
  the $i$-th root tag answer is sampled;
\item no candidate's $\ROflat$ value is read before its append step, so
  the sequence satisfies the predictability and unrevealed-value premises
  of \Cref{lem:adaptive-candidate};
  and
\item its length is at most $\qRO(\mathcal{B}) + m$.
\end{enumerate}
\end{lemma}

\begin{proof}
The committing, discoverability, and hash (collision, second-preimage,
and preimage) games are keyless: all keys are
adversary-supplied, so a direct query to a keyed transcript is one of
$\mathcal{B}$'s visible $\ROflat$ queries counted in $\qRO(\mathcal{B})$,
and there is no $\badkey$ event. No construction oracle runs while
$\mathcal{B}$ issues its direct queries, so each direct-query candidate is
first read no earlier than its append step. The append predicate depends on $Y$ alone
through the exact decoder of \Cref{lem:bridge-decode}, which ignores the
requested output length, and by Output-Point Separation
(\Cref{cor:output-sep}) each input decodes to at most one output-point
class; hence each distinct direct query contributes at most one
candidate, giving~(1).

If the challenger reaches its $m$ checks, append $\flatmap(R_i)$
immediately before the $i$-th root tag answer is sampled, unless it is
already present. The intermediate $\ROflat$ lookups of check $i$ lie at
non-root tag output points, so by \Cref{cor:output-sep} they do not read
$\flatmap(R_i)$. If $\flatmap(R_i)$ was touched earlier---by a direct
query or by an earlier check $j < i$---then \Cref{lem:bridge-decode} and
\Cref{cor:output-sep} make that earlier access a final-root lookup at the
same address, which was therefore already appended, so the input is in
the sequence and is not appended again. Otherwise its value is still
unread when appended. Either way no candidate's value is read before its
append step, giving~(2) and~(3). At most $\qRO(\mathcal{B})$ direct
candidates and $m$ challenger candidates are appended, giving~(4).
\end{proof}

\begin{lemma}[Root Tag Hash Bound]
\label{lem:root-tag-hash}
Consider a keyless \TW{} random oracle game (\Cref{sec:ro-model}) in
which, after $\mathcal{B}$'s direct $\ROflat$ queries, the challenger runs
deterministic encryption or decryption checks reaching final root
transcripts, with win event $\mathsf{W}$. By the Root tag Candidate
Sequence lemma (\Cref{lem:root-candidate-seq}) the root tag inputs form a
chronological candidate sequence meeting the premises of
\Cref{lem:adaptive-candidate} with $\tau = \tauroot$. Writing
$\qRO = \qRO(\mathcal{B})$:
\begin{enumerate}
\item if $\mathsf{W}$ forces $s$ checks to reach pairwise distinct bridged
  final root transcripts carrying a common root tag, then, with $m = s$
  checks, $\Pr[\mathsf{W}] \leq \RootColl_{\tauroot}(\qRO, s)$;
\item if $\mathsf{W}$ forces a single check ($m = 1$) to carry a root tag
  equal to a target fixed independently of $\ROflat$, then
  $\Pr[\mathsf{W}] \leq \RootPre_{\tauroot}(\qRO)$;
\item if a designated check reaches $R_X$ whose root tag is exposed before
  $\mathcal{B}$ outputs, and $\mathsf{W}$ forces a single further check
  ($m = 1$ besides $R_X$) to reach a transcript $R' \neq R_X$ carrying
  $R_X$'s root tag, then $\Pr[\mathsf{W}] \leq \RootPre_{\tauroot}(\qRO)$.
\end{enumerate}
\end{lemma}

\begin{proof}
In each case the candidate sequence has the stated length and meets the
premises of \Cref{lem:adaptive-candidate}.
\emph{(1)} The $s$ distinct candidates carrying a common
$\tauroot$-projection are an $s$-way collision; the collision case with
$L = \qRO + s$ gives
$\binom{\qRO + s}{s} \cdot 2^{-\tauroot(s-1)}
= \RootColl_{\tauroot}(\qRO, s)$.
\emph{(2)} The candidate hitting the fixed target is a preimage; the
preimage case with $L = \qRO + 1$ gives
$(\qRO + 1) \cdot 2^{-\tauroot} = \RootPre_{\tauroot}(\qRO)$.
\emph{(3)} Prepend $\flatmap(R_X)$ as the designated candidate, its value
the exposed root tag; the further candidate colliding with it on its
$\tauroot$-projection is a second-preimage, and the second-preimage
corollary with $L = \qRO + 2$ gives
$(\qRO + 1) \cdot 2^{-\tauroot} = \RootPre_{\tauroot}(\qRO)$.
\end{proof}

\begin{theorem}[Random Oracle CMTDs-4]
\label{thm:cmtds4-ro}
For every $s \geq 2$ and every $s$-way CMTDs-4
adversary $\mathcal{B}$ in the \TW{} random oracle model
(\Cref{sec:ro-model}),
\[
  \Advcmtds{4}_\RO(\mathcal{B})
  \;\leq\; \RootColl_{\tauroot}(\qRO(\mathcal{B}), s)
  = \frac{\binom{\qRO(\mathcal{B}) + s}{s}}{2^{\tauroot (s - 1)}}.
\]
\end{theorem}

\begin{proof}
Run the random oracle CMTDs-4 game. After rejecting on any failed syntax,
message-length, or pairwise-full-input check, the challenger performs $s$
deterministic decryption checks $\Dec_{K_i}^\RO(N_i,A_i,C \concat T)$ on
the one common body $C \concat T$, reaching $R_1,\dots,R_s$. On a winning
execution the tuples $(K_i,N_i,A_i,M_i)$ are pairwise distinct, and since
all checks decrypt the common body, determinism of $\Dec$ makes the
triples $(K_i,N_i,A_i)$ pairwise distinct---equal triples would return
equal messages, hence equal tuples. By Opening Triple Separation
(\Cref{cor:triple-sep}) the candidates $\flatmap(R_1),\dots,\flatmap(R_s)$
are pairwise distinct (leaf tag collisions cannot merge them), and every
accepting check carries the common root tag $T$. The Root Tag Hash Bound
(\Cref{lem:root-tag-hash}, part~1) gives
$\Advcmtds{4}_\RO(\mathcal{B}) \leq \RootColl_{\tauroot}(\qRO(\mathcal{B}), s)$.
\end{proof}

\paragraph{Remaining committing notions.}
The hierarchy and $s$-way extension of the \cite{BH22} committing notions
are given in \cite[\S3, Figs.~4--5, and App.~A]{BH22}. We use those
relations for the D-notions and a direct source-game argument for the
CMTs-$\ell$ notions to preserve the resource split of
\Cref{sec:resources}.

\begin{corollary}[Committing-Notion Root-Collision Bound]
\label{cor:committing-combined}
For every $s \geq 2$, every
\[
  X \in \bigl\{\mathsf{CMTDs\text{-}1},\, \mathsf{CMTDs\text{-}3},\,
              \mathsf{CMTDs\text{-}4},\, \mathsf{CMTs\text{-}1},\,
              \mathsf{CMTs\text{-}3},\, \mathsf{CMTs\text{-}4}\bigr\},
\]
and every $s$-way $X$-adversary $\mathcal{A}$ against \TW{},
\[
  \Adv^X_\TW(\mathcal{A})
  \;\leq\;
  \Lift_c(\Ntot) + \RootColl_{\tauroot}(\Nprim, s)
  \;=\;
  \Lift_c(\Ntot) + \frac{\binom{\Nprim + s}{s}}{2^{\tauroot(s-1)}}.
\]
\end{corollary}

\begin{proof}
CMTDs-4 is bounded directly by \Cref{cor:cmtds4}. For
$\ell \in \{1,3\}$, the $s$-way hierarchy of \cite[\S3, Fig.~5 and
App.~A]{BH22} gives
$\Advcmtds{\ell}_\TW(\mathcal{A}) \leq
\Advcmtds{4}_\TW(\mathcal{A})$ for the same output distribution, so the
resource counts are unchanged.

The E-notions are bounded by the Root Tag Hash Bound, with the
challenger's $s$ deterministic \emph{encryption} checks supplying the
candidate final-root transcripts
(\Cref{lem:root-candidate-seq} builds the candidate sequence for
encryption or decryption checks alike). On a winning execution all $s$
encryptions are equal, hence of equal length and with one common root tag
$T$; pairwise $\WiC_\ell$-distinctness forces pairwise distinct
triples---for $\ell = 4$ via determinism, since equal triples with equal
ciphertexts would decrypt to equal messages---so \Cref{cor:triple-sep}
makes the $s$ final-root candidates distinct while their $\tauroot$-bit
projections all equal $T$. By the Lift Application corollary (\Cref{cor:lift-application}) for the
$X$ game, the derived random oracle adversary $\mathcal{B}$ has
$\qRO(\mathcal{B}) \leq \Nprim$; the Root Tag Hash Bound
(\Cref{lem:root-tag-hash}, part~1) then bounds its win by
$\RootColl_{\tauroot}(\qRO(\mathcal{B}), s)$ in the random oracle model,
which the lift transfers to the public permutation setting at $\Nprim$;
the encryption checks are construction-internal work counted in $N$, not
in $\Nprim$.
\end{proof}

\begin{theorem}[Random Oracle Collision Resistance of $\HTW$]
\label{thm:coll-ro}
For every $s \geq 2$ and every $s$-way collision
adversary $\mathcal{B}$ in the \TW{} random oracle model,
\[
  \AdvColl_{\RO}(\mathcal{B})
  \;\leq\;
  \LeafColl_{\tauleaf}(\qRO(\mathcal{B})+\nleaf(\mathcal{B}))
  + \RootColl_{\tauroot}(\qRO(\mathcal{B}), s).
\]
\end{theorem}

\begin{proof}
Run the random oracle collision game. Let $\badleaf^\star$ be the event
that two distinct construction-generated final leaf transcripts receive
the same $\tauleaf$-bit $\RF$ projection; by the adversary-supplied-key
bound of \Cref{lem:leaf-collision},
$\Pr[\badleaf^\star] \leq \LeafColl_{\tauleaf}(\qRO(\mathcal{B})+\nleaf(\mathcal{B}))$,
so it remains to bound the win under $\neg\badleaf^\star$. After rejecting
unless the $s$ output inputs are pairwise distinct and in the \TW{}
domain, the challenger performs $s$ deterministic encryptions
$\Enc_{K_i}^{\RO}(N_i,A_i,M_i)$, reaching $R_1,\dots,R_s$, all carrying a
common root tag $T$ on a win. Pairwise distinct tuples $(K_i,N_i,A_i,M_i)$
have pairwise distinct $(K_i,N_i,A_i,C_i)$: two tuples differing in their
context already differ in $(K,N,A)$, while two sharing a context but
differing in the message yield distinct ciphertexts, since for a fixed
context encryption maps the message to the ciphertext bijectively. So on
$\neg\badleaf^\star$ Final-Root Input Separation
(\Cref{prop:final-root-sep}) makes $\flatmap(R_1),\dots,\flatmap(R_s)$
pairwise distinct. The Root Tag Hash Bound (\Cref{lem:root-tag-hash},
part~1) bounds the win under $\neg\badleaf^\star$ by
$\RootColl_{\tauroot}(\qRO(\mathcal{B}),s)$; adding $\Pr[\badleaf^\star]$
gives the theorem.
\end{proof}

\begin{theorem}[Random Oracle Everywhere Preimage Resistance of $\HTW$]
\label{thm:epre-ro}
For every query-independent target tag $T^* \in \bits^{\tauroot}$ and
every preimage adversary $\mathcal{B}$ in the \TW{}
random oracle model (\Cref{sec:ro-model}),
\[
  \AdvePre_\RO(\mathcal{B})
  \;\leq\; \RootPre_{\tauroot}(\qRO(\mathcal{B}))
  = \frac{\qRO(\mathcal{B}) + 1}{2^{\tauroot}}.
\]
\end{theorem}

\begin{proof}
Run the random oracle preimage game for the query-independent target
$T^*$, which is independent of $\ROflat$. After rejecting unless the
output input $X = (K,N,A,M)$ is in the \TW{} domain, the challenger
performs one encryption $\Enc_{K}^\RO(N,A,M)$, reaching a final root
transcript $R$. A win means $\HTW(X) = T^*$, i.e.\ the root tag
$\lfloor \ROflat(\flatmap(R)) \rfloor_{\tauroot}$ at $R$ equals the fixed
target $T^*$. The Root Tag Hash Bound (\Cref{lem:root-tag-hash}, part~2,
with target $T^*$) gives
$\AdvePre_\RO(\mathcal{B}) \leq \RootPre_{\tauroot}(\qRO(\mathcal{B}))$.
\end{proof}

\paragraph{Context discoverability.}

\begin{proof}[Proof of \Cref{thm:cdy-full}]
Full reveal $\mathsf{ts} = \{k,n,a\}$ hides no component and is won
trivially by any scheme, so the restriction to
$\mathsf{ts} \subsetneq \{k,n,a\}$ excludes no meaningful case. The
restricted advantage requires $\mathcal{A}$ to win against every
query-independent selector, so it is at most the success probability
against the selector $S^\star$ that samples a fresh context
$(K_0, N_0, A_0)$---each of $K_0$, $N_0$, and $A_0$ uniform on a
$256$-bit space---encrypts a message under it to obtain the target
ciphertext $C^* = C \concat T^*$, and reveals only the
$\mathsf{ts}$-components. By the Lift Application corollary
(\Cref{cor:lift-application}), the derived random oracle adversary
$\mathcal{B}$ has $\qRO(\mathcal{B}) \leq \Nprim$; it remains to bound the
random oracle success against $S^\star$. Here
$T^* = \lfloor \RF(R_0) \rfloor_{\tauroot}$ is the root tag of the
selector's final root transcript $R_0$ under $(K_0, N_0, A_0)$, and a
winning $\mathcal{A}$ outputs a context reaching a final root transcript
$R$ with $\RF(R)$ projecting to $T^*$.

Let $Z_0$ be any context component hidden by $\mathsf{ts}$---one exists
because $\mathsf{ts} \subsetneq \{k,n,a\}$---which $\mathcal{A}$ must itself
supply (for the named variants, the key for
$\mathsf{ts} \in \{\varnothing, \{n,a\}\}$, the nonce for
$\mathsf{ts} = \{k,a\}$, the associated data for $\mathsf{ts} = \{k,n\}$),
and which $S^\star$ sampled uniformly on $\bits^w$
with $w = 256 \geq k$ and which $\mathcal{A}$ never sees: the only values
revealed to it, the body and tag of $C^*$, are $\RF$ projections
independent of $Z_0$. If $R = R_0$, then $\mathcal{A}$ has reconstructed
the selector's transcript, in particular its $Z_0$-component, which by
\Cref{lem:bridge-decode} requires a candidate decoding under
$\KeyedTWOut_{Z_0}$ to the hidden $Z_0$; since $Z_0$ stays uniform over
the values no direct query has tested, \Cref{lem:key-tests} (hidden
component $Z_0$, $w = 256$) bounds this by $\KeyGuess_k(\Nprim)$. If
$R \neq R_0$, the win is that $R$ collides with $R_0$ on its
$\tauroot$-bit projection, which the Root tag Candidate Sequence lemma
(\Cref{lem:root-candidate-seq}) and the second-preimage corollary of
\Cref{lem:adaptive-candidate}, with the designated candidate
$\flatmap(R_0)$, bound by $\RootPre_{\tauroot}(\Nprim)$. Adding the lift
gives the bound, uniformly over $\mathsf{ts}$.
\end{proof}

\begin{theorem}[Random Oracle aSec]
\label{thm:asec-ro}
For every aSec adversary $\mathcal{B}$ in the \TW{} random oracle model
(\Cref{sec:ro-model}),
\[
  \AdvaSec_\RO(\mathcal{B})
  \;\leq\;
  \RootPre_{\tauroot}(\qRO(\mathcal{B}))
  + \LeafColl_{\tauleaf}(\qRO(\mathcal{B})+\nleaf(\mathcal{B})).
\]
\end{theorem}

\begin{proof}
Run the random oracle aSec game: the challenger draws the challenge input
$X$, encrypts it to obtain $T = \HTW(X)$ at a final root transcript $R_X$,
and gives $X$ to $\mathcal{B}$, which outputs $X' \neq X$, winning if
$\HTW(X') = T$. Let $\badleaf^\star$ be the event that two distinct
construction-generated final leaf transcripts receive the same
$\tauleaf$-bit $\RF$ projection; by \Cref{lem:leaf-collision},
$\Pr[\badleaf^\star] \leq \LeafColl_{\tauleaf}(\qRO(\mathcal{B})+\nleaf(\mathcal{B}))$,
so it remains to bound the win under $\neg\badleaf^\star$. Encrypting $X'$
reaches a final root transcript $R'$. If $\flatmap(R') = \flatmap(R_X)$
then on $\neg\badleaf^\star$ Final-Root Input Separation
(\Cref{prop:final-root-sep}) makes $X$ and $X'$ share $(K,N,A,C)$, hence
the same message by determinism, so $X' = X$---a contradiction; thus
$R' \neq R_X$. The challenge root tag
$T = \lfloor \RF(R_X) \rfloor_{\tauroot}$ is exposed to $\mathcal{B}$, and
a win places $R'$'s root tag at $T$. The Root Tag Hash Bound
(\Cref{lem:root-tag-hash}, part~3, with designated $R_X$) bounds the win
under $\neg\badleaf^\star$ by $\RootPre_{\tauroot}(\qRO(\mathcal{B}))$;
adding $\Pr[\badleaf^\star]$ gives the theorem.
\end{proof}
